% =============================================================================
% SECTION 2: MODEL
% =============================================================================

\section{The Model}

This section presents a formal model of firm compliance decisions in a compute permit market with imperfect monitoring. We derive the conditions under which firms choose to comply with permit requirements and characterize how regulatory parameters affect aggregate compliance.

\subsection{Setup}

Consider a market with $N$ firms, indexed by $i \in \{1, \ldots, N\}$. Each firm seeks to conduct a training run requiring compute $C_i$, measured in FLOP. A regulatory threshold $\tau$ (e.g., $10^{25}$ FLOP) determines which training runs require permits. Firms with $C_i > \tau$ must hold a valid permit to operate legally.

Each firm has a private valuation $v_i$ representing the economic benefit of completing its training run. This valuation captures expected revenues from deploying the trained model, strategic advantages from technological leadership, and other sources of value. Firms are heterogeneous in their valuations, reflecting differences in business models, market positions, and technical capabilities.

Firms face a binary choice: comply with permit requirements or defect by operating without a permit. A compliant firm purchases the required permits at market price $P$ and operates legally. A defecting firm avoids permit costs but faces the risk of detection and penalty. Throughout, monetary values are expressed in millions of USD (M\$) and compute is measured in FLOP.

\subsection{The Deterrence Condition}

A firm's decision to comply depends on comparing the cost of compliance against the expected cost of defection. Let:

\begin{itemize}
    \item $g_i$ denote the gain from defection for firm $i$.
    \item $p_{\mathrm{detect},i}$ denote the effective probability of detecting firm $i$'s violation.
    \item $B_i$ denote the total perceived penalty upon detection.
\end{itemize}

A risk-neutral firm complies if and only if the expected penalty from defection exceeds the gain:

\begin{equation}
p_{\mathrm{detect},i} \times B_i \geq g_i
\label{eq:deterrence}
\end{equation}

This is the central deterrence condition of the model. When it holds, compliance is the economically rational choice. When it fails, firms have an incentive to defect.

The gain from defection for firm $i$ is:

\begin{equation}
g_i = \min(p^*, v_i) + r_i \, V_i
\label{eq:gain}
\end{equation}

where $p^*$ is the current market clearing price for permits, $v_i$ is the firm's private valuation, $V_i$ is the capability value of running a frontier model, and $r_i$ is a firm-specific racing factor capturing competitive urgency. The $\min$ operator reflects that a firm whose valuation falls below the permit price would not have purchased a permit regardless, so the savings from evasion are bounded by $v_i$. In most calibrations $r_i \, V_i = 0$, so the racing term drops out and the gain reduces to the lesser of the permit price and the firm's valuation.

\subsection{The Permit Market}
\label{sec:permit-market}

The clearing price $p^*$ that enters the gain from defection is determined by a permit market. A permit grants the right to conduct compute up to a threshold $F_{\text{permit}}$ FLOP---it is a \emph{quantity} licence, not a binary permission. A firm requiring a run of size $C_i$ submits a demand of $\lceil C_i / F_{\text{permit}} \rceil$ permit units based on its declared compute needs. Under the baseline calibration, $F_{\text{permit}}$ is set equal to one standard frontier training run, so the large majority of firms demand either zero or one permit and the mechanics appear nearly binary---but the general case accommodates multi-permit demand for larger runs or smaller permit units. The permit supply is fixed at $Q$ units per period, and the market clears each period against this cap.

The model supports two clearing modes, chosen per scenario to reflect different regulatory philosophies. Auction mode allows price to emerge from competition, making the permit cost endogenous to compliance behaviour. Fixed-price mode decouples price from demand, letting the regulator directly set the permit-cost term that enters each firm's deterrence calculation---at the cost of sacrificing the informational content of market prices.

In \emph{auction mode}, each firm $i$ submits a triple $(\text{id}_i,\, q_i,\, b_i)$, where $q_i$ is the number of permits demanded and $b_i$ is the willingness-to-pay per permit. Bids are expanded into $q_i$ individual permit-unit offers each priced at $b_i$, sorted descending, and the top $Q$ units are allocated. The market-clearing price $p^*$ is the bid of the marginal ($Q$th highest) unit; all winning firms pay this uniform price. When aggregate demand falls weakly below supply---$\sum_i q_i \leq Q$---every firm receives its requested allocation at $p^* = 0$. This zero-price outcome is deliberate: a market with slack imposes no scarcity rent. It does, however, imply that under universal access ($Q = N$) an auction would always clear at $p^* = 0$, and the permit-cost term would drop out of the gain in Equation~\eqref{eq:gain}. This is why the universal-access and enforcement-design scenarios adopt fixed-price mode, where the regulator can post a non-zero permit cost regardless of market slack.

In \emph{fixed-price mode}, the regulator posts an administrative price $\bar{p}$ rather than discovering price through competition. Any firm whose bid $b_i \geq \bar{p}$ qualifies for permits at cost $\bar{p}$ per unit; if qualifying demand exceeds $Q$, permits are allocated by uniform random draw among qualifying units up to the cap.\footnote{Under all reported calibrations the permit cap equals the firm count ($Q = N = 20$) and permit demand is binary, so the number of qualifying bidders never exceeds supply and the random rationing rule is never invoked; it is retained for generality and for fixed-price configurations with $Q < N$.} Firms that do not win the draw obtain no permit and face the deterrence decision as evaders, with their avoided permit cost evaluated at the posted price; that is, $p^* = \bar{p}$ enters their gain in Equation~\eqref{eq:gain}. Firms that do not win the draw obtain no permit and face the deterrence decision as evaders, with their avoided permit cost evaluated at the posted price; that is, $p^* = \bar{p}$ enters their gain in Equation~\eqref{eq:gain}.

A potential consequence of the auction case is a compliance externality: if non-compliant firms were to abstain from permit bidding, aggregate demand would soften and $p^*$ would fall---reducing the cost of compliance avoided by evasion, weakening the deterrence condition for marginal evaders, and creating a self-reinforcing dynamic toward a low-compliance equilibrium. In the simulations reported here this channel is held closed: every firm above the threshold submits its bid each period regardless of its subsequent compliance choice, so demand---and hence $p^*$---does not respond to the compliance rate. Fixed-price mode eliminates the channel by construction, since the permit cost is invariant to bidding behaviour, which is one reason the enforcement-design scenarios adopt it. Endogenizing bid abstention is a natural extension for studying the stability of auction-based permit regimes.

\subsection{Detection Probability}
\label{sec:detection}

Detection follows a two-stage sequential process. In the first stage, firm $i$ is selected for audit with probability $p_{\mathrm{audit}}(i)$. In the second stage, conditional on an audit occurring, the violation is detected with probability $p_{\mathrm{catch}}$. The overall detection probability is therefore:

\begin{equation}
p_{\mathrm{detect}}(i) = p_{\mathrm{audit}}(i) \times p_{\mathrm{catch}}
\label{eq:detect}
\end{equation}

The audit is the sole mechanism that can produce a legally demonstrable finding. Whistleblower reports, hardware monitoring, and historical backchecks serve as information sources that improve the probability of detection \emph{within} an audit, rather than as independent enforcement channels. This reflects the institutional reality that evidence gathered outside a formal audit process typically cannot, on its own, establish a regulatory violation.

The audit probability depends on the regulatory mode. Under random auditing, all firms face the base audit rate:

\begin{equation}
p_{\mathrm{audit}}(i) = \pi_0
\label{eq:audit-random}
\end{equation}

Under signal-dependent targeting, the audit probability also reflects the strength of the suspicion signal generated by the firm's unpermitted excess compute:

\begin{equation}
p_{\mathrm{audit}}(i) = \min\!\Big(1,\; \pi_0 + c(i) \cdot s(i) \cdot (1 - \pi_0)\Big)
\label{eq:audit-signal}
\end{equation}

where $c(i) \geq 0$ is a firm-specific audit coefficient reflecting visibility or regulatory scrutiny, and $s(i) \in [0,1]$ is the suspicion signal, defined below.

Firms whose planned training run exceeds the regulatory threshold $\tau$ generate a suspicion signal proportional to their unpermitted excess compute:

\begin{equation}
s(i) = \min\!\left(1,\ \left(\frac{e(i)}{\tau}\right)^{\!\gamma}\right)
\label{eq:signal}
\end{equation}

where $e(i)$ is firm $i$'s unpermitted excess FLOP and $\gamma > 0$ is the signal exponent. At $\gamma = 1$ the mapping is linear; $\gamma < 1$ makes small violations more visible (concave), while $\gamma > 1$ makes only large violations prominent (convex). Compliant firms carry no unpermitted excess compute, so $s(i) = 0$ and their audit probability collapses to the random floor $\pi_0$. The baseline rate $\pi_0$ thus represents the minimum audit exposure guaranteed by the regulatory regime---akin to random compliance checks---while $c(i)$ scales how much a firm's excess compute raises its audit probability above that floor. Note that at $s(i) = 1$ and $c(i) = 1$, equation~\eqref{eq:audit-signal} yields $p_{\mathrm{audit}} = 1$: the regulator audits firms whose excess compute reaches the threshold with certainty. If a maximum number of audits per period $N_{\max}$ is set, the regulator prioritises: under signal-dependent auditing, triggered firms are ranked by $s(i)$ descending and the top $N_{\max}$ are selected; under random auditing, $N_{\max}$ triggered firms are sampled uniformly.

Conditional on an audit, detection proceeds through four nested channels in sequence: a direct audit pass, a historical backcheck, whistleblower disclosure, and hardware monitoring. The violation is caught if any channel fires:

\begin{equation}
p_{\mathrm{catch}} = 1 - \beta \cdot (1 - p_b) \cdot (1 - p_w) \cdot (1 - p_m)
\label{eq:catch}
\end{equation}

where $\beta \in [0,1]$ is the false negative rate of the direct audit pass, $p_b$ is the backcheck probability, $p_w$ is the whistleblower detection probability, and $p_m$ is the monitoring detection probability. The miss probability is the joint probability that all four channels fail. This sequential nesting also allows each detection event to be attributed to a single channel in the simulation output.

Audits are imperfect in both directions. Compliant firms face a false positive rate $\alpha$ on the direct pass and can be flagged by the backcheck channel; whistleblower and monitoring channels do not generate false positives, as there is no real violation to find. False positives impose administrative costs on compliant firms and consume regulatory resources without producing enforcement value, a burden quantified in Section~\ref{sec:results-workload}.

\subsection{Penalty Structure}

The total perceived penalty $B_i$ upon detection is firm-specific:

\begin{equation}
B_i = (\phi + K + R_i) \times \rho_i
\label{eq:burden}
\end{equation}

where $\phi$ is the monetary fine imposed upon violation, $K$ is the collateral forfeited, $R_i$ is firm $i$'s perceived reputational cost if caught, and $\rho_i \geq 0$ is a risk profile multiplier scaling the firm's sensitivity to losses. Collateral is posted upfront before the firm operates and is returned if no violation is detected. This structure addresses the limited liability problem: firms with insufficient assets to pay fines can still be deterred if they must post collateral in advance. Financial penalties are bounded by the firm's capacity to pay: the effective monetary sanction is $P_{\mathrm{eff}} = \min(K + \phi,\, L)$, where $L$ is the liability cap. Because collateral is collected before the violation occurs, it restores deterrence even when the ex-post fine alone is insufficient. The reputation term $R_i$ introduces heterogeneity in how firms perceive the consequences of detection beyond direct financial penalties, and is the mechanism through which repeated violations compound deterrence in dynamic scenarios (see Appendix~D).

\subsection{Dynamic Enforcement}
\label{sec:dynamics}

The audit coefficient and reputation sensitivity introduced above are static in the baseline model. When dynamic updating is enabled, both evolve in response to enforcement outcomes, allowing the model to capture feedback between a firm's compliance history and the regulatory pressure it faces. Following a detected violation, the firm's audit coefficient increases by an escalation increment; the excess above the baseline then decays exponentially toward the base at the end of each period, with the escalation applied before the within-period decay:

\begin{equation}
c_i^{(t)} = c_i^{\mathrm{base}} + \Bigl(c_i^{(t-1)} - c_i^{\mathrm{base}} + \Delta \cdot \mathbf{1}[\text{caught}_{t}]\Bigr)\cdot(1 - \delta)
\label{eq:coeff}
\end{equation}

where $c_i^{\mathrm{base}}$ is the firm's configured baseline coefficient (the long-run floor), $\delta \in [0,1]$ is the per-period decay rate, and $\Delta \geq 0$ is the escalation increment per failure. Each failed audit also multiplies the firm's reputation sensitivity:

\begin{equation}
R_i^{(t)} = R_i^{(0)} \cdot \bigl(1 + \varepsilon\bigr)^{n_i}
\label{eq:rep}
\end{equation}

where $\varepsilon \geq 0$ is the reputation escalation factor and $n_i$ is the firm's cumulative count of failed audits. Reputation escalation is permanent (no decay); setting $\Delta = \varepsilon = 0$ recovers a static enforcement environment.

These two mechanisms were chosen as representative of two forms of feedback-driven deterrence: strong feedback (exponential increase) and weak feedback (capped increase with decay). Their per-period effects on compliance trajectories are explored empirically in Appendix~\ref{sec:results-dynamic}.

\subsection{Equilibrium Compliance}

Given the deterrence condition, aggregate compliance has two components. Firms that require no additional permits comply automatically: this includes firms whose planned compute falls below the regulatory threshold $\tau$ (no permit is required) and firms holding enough permits to cover their full compute (the purchase itself constitutes compliance). Let $\mathcal{P}$ denote this set. Only firms with unpermitted excess compute evaluate the deterrence condition. The compliance rate is therefore:

\begin{equation}
\text{Compliance Rate} = \frac{|\mathcal{P}|}{N} + \frac{|\mathcal{U}|}{N} \cdot \Pr\!\big(p_{\mathrm{detect},i} \times B_i \geq g_i \;\big|\; i \in \mathcal{U}\big)
\end{equation}

where $\mathcal{U}$ is the set of firms with unpermitted excess compute. In supply-constrained scenarios where $Q < N$, the permitted portion of the first term is bounded by $Q/N$---plus the fraction of firms below the threshold---and the second term captures any additional compliance induced by enforcement. In universal-access scenarios where $Q = N$ and all firms can obtain permits, the distinction collapses and the compliance rate depends entirely on whether permit purchase is individually rational.

This expression clarifies the levers available to regulators. Compliance can be increased by:

\begin{enumerate}
    \item Increasing $p_{\mathrm{detect},i}$ through higher audit rates, better monitoring, or whistleblower programs.
    \item Increasing $B_i$ through higher fines, collateral requirements, or reputation costs.
    \item Decreasing $g_i$ through more efficient permit markets that reduce the cost of compliance.
\end{enumerate}

The smart enforcement approach emphasizes the third lever: by making permits cheaper and more accessible, regulators reduce the gain from defection, allowing compliance to be maintained even with moderate detection probabilities.

\subsection{Model Parameters}

Table \ref{tab:parameters} summarizes the key parameters of the model. Detailed calibration values for each regulatory scenario are provided in Appendix A.

\begin{table}[h]
\centering
\caption{Key Model Parameters}
\label{tab:parameters}
\begin{tabular}{@{}llp{7cm}@{}}
\toprule
\textbf{Parameter} & \textbf{Symbol} & \textbf{Description} \\
\midrule
Base audit rate & $\pi_0$ & Probability any firm is audited each period \\
Audit coefficient & $c(i)$ & Firm-specific signal scaling (signal mode only) \\
Signal exponent & $\gamma$ & Shape of the excess-compute $\to$ signal curve \\
False negative rate & $\beta$ & Probability direct audit pass misses a violation \\
False positive rate & $\alpha$ & Rate of spurious findings on compliant firms \\
Backcheck probability & $p_b$ & Historical review catch rate (within audit) \\
Whistleblower probability & $p_w$ & Disclosure-based catch rate (within audit) \\
Monitoring probability & $p_m$ & Hardware/metering catch rate (within audit) \\
Audit capacity & $N_{\max}$ & Maximum audits conducted per period \\
Penalty amount & $\phi$ & Fine imposed upon detected violation \\
Collateral & $K$ & Upfront deposit forfeited upon violation \\
Liability cap & $L$ & Upper bound on total monetary sanction \\
Reputation sensitivity & $R_i$ & Perceived brand/trust cost if caught \\
Risk profile & $\rho_i$ & Multiplier on perceived losses \\
Capability value & $V_i$ & Strategic value of running a frontier model \\
Racing factor & $r_i$ & Urgency multiplier on capability value \\
FLOP threshold & $\tau$ & Compute level above which permits are required \\
Permit size & $F_{\text{permit}}$ & Compute entitlement per permit unit \\
Permit supply & $Q$ & Permit units available per period \\
Fixed permit price & $\bar{p}$ & Administrative price (fixed-price mode only) \\
Coefficient decay & $\delta$ & Per-period decay of audit coefficient (dynamic mode) \\
Escalation increment & $\Delta$ & Audit-coefficient increase per detection (dynamic mode) \\
Reputation escalation & $\varepsilon$ & Reputation multiplier per failed audit (dynamic mode) \\
\bottomrule
\end{tabular}
\end{table}
